\chapter{Dérivations mathématiques}
\label{app:derivations}

Cette annexe rassemble les justifications algébriques des formules employées dans
le corps du document. Les chapitres principaux en gardent les expressions et leur
interprétation opérationnelle.

\section{Inversion du logit}
\label{app:logit-inversion}

En posant $p=\Pp(Y=1\mid X=\x)$ et $z=\thv^\top\xt$,
\begin{align*}
  \log\frac{p}{1-p}=z
  &\iff \frac{p}{1-p}=e^z
  \iff p=e^z(1-p)\\
  &\iff p(1+e^z)=e^z
  \iff p=\frac{e^z}{1+e^z}=\frac1{1+e^{-z}}.
\end{align*}

\section{Gradient de la log-perte}
\label{app:gradient-derivation}

Pour une observation,
$\ell_i=-y_i\log p_i-(1-y_i)\log(1-p_i)$ et
$\partial p_i/\partial z_i=p_i(1-p_i)$. La règle de chaîne donne
\begin{align*}
  \frac{\partial\ell_i}{\partial z_i}
  &=-\frac{y_i}{p_i}p_i(1-p_i)
    +\frac{1-y_i}{1-p_i}p_i(1-p_i)\\
  &=p_i-y_i.
\end{align*}
Comme $\nabla_{\thv}z_i=\xt_i$, la moyenne des contributions vaut
\[
  \nabla J(\thv)=\frac1m\sum_{i=1}^m(p_i-y_i)\xt_i
  =\frac1m\X^\top\bigl(\sigma(\X\thv)-\y\bigr).
\]

\section{Convexité de la log-perte}
\label{app:convexity-proof}

Avec $W=\operatorname{diag}(p_i(1-p_i))$,
\[
  \nabla^2J(\thv)=\frac1m\X^\top W\X.
\]
Pour tout $v\in\R^{n+1}$,
\[
  v^\top\nabla^2J(\thv)v
  =\frac1m(\X v)^\top W(\X v)
  =\frac1m\sum_{i=1}^m p_i(1-p_i)(\xt_i^\top v)^2\geq0.
\]
La Hessienne est donc positive semi-définie et $J$ est convexe.

\section{Direction de descente}
\label{app:descent-direction}

Pour un déplacement suffisamment petit $\Delta\thv$,
\[
  J(\thv+\Delta\thv)
  =J(\thv)+\nabla J(\thv)^\top\Delta\thv
  +o(\|\Delta\thv\|).
\]
Avec $\Delta\thv=-\alpha\nabla J(\thv)$,
\[
  J(\thv+\Delta\thv)
  =J(\thv)-\alpha\|\nabla J(\thv)\|_2^2
  +o(\alpha\|\nabla J(\thv)\|).
\]
Le terme principal est négatif pour $\alpha>0$ lorsque le gradient est non nul.
